% !TEX root = ../main.tex
\section{Conclusion}
\label{sec:conclusion}

We presented \medskille, a framework for evaluating and evolving medical agent skills at scale.
\medskillbench provides the first skill-level execution benchmark for medical AI, covering 1,462 skills with eight-dimensional scoring in isolated sandbox environments.
Using this benchmark, we identified and quantified a systematic \emph{read-but-not-use} gap of 22 percentage points between skill invocation (96.2\%) and task completion (74.3\%).
Critically, stratification by \texttt{tool\_type} reveals that guide-style skills without executable interfaces exhibit the largest gap (0.269), while tool-based skills achieve substantially tighter coupling between comprehension and execution (0.069--0.111).
This finding reframes the core bottleneck: the primary failure mode is not that models cannot operate tools, but that natural-language-only skill descriptions lack sufficient grounding for reliable execution.

Building on this diagnostic signal, we proposed \gepamed, a training-free method that uses execution trajectories and judge feedback to iteratively refine skill descriptions.
\todo{Report GEPA-Med improvement magnitude and cross-model transfer results once experiments conclude.}

Our results establish that medical agent skills are not static documents to be written once and deployed indefinitely.
They are evolvable interfaces whose quality can be systematically measured and improved through evaluation-driven feedback loops.
We believe the evaluate--diagnose--evolve paradigm extends naturally beyond the medical domain to any setting where LLM agents consume structured skill specifications, offering a training-free path toward more reliable agent--tool interaction.
We discuss remaining limitations in \S\ref{sec:limitations}.
